% secciones/01_introduccion/introduccion.tex
\section{Introducción}

En los ultimos años se ha visto un aumento en el acceso a peliculas a traves de plataformas de streaming, estas plataformas cuentan con un gran numero de usuarios los cuales en muchas ocasiones tienen la necesidad de poder contar con una herramienta que les permita encontrar contenido que se ajuste a sus intereses.

En este proyecto se propone el desarollo de un asistente virtual para la recomendacion de peliculas el cual permitira a los usuarios realizar sus preguntas en español sobre las peliculas del catalogo y recibir informacion sobre las mismas, teniendo en cuenta caracteristias como titulo, genero, sinopsis, entre otros.

Es en este contexto que el poder implementar un sistema de clasificación automatica de texto resulta realmente util, el procesamiento PLN (Procesamiento del Lenguaje Natural) permite desarrollar sistemas capaces de procesar y analizar textos dados por el usuario y relacionarlos directamente con la informacion disponible, en este caso, informacion sobre peliculas.

Existen diferentes tecnicas que pueden considerarse para poder implementar un sistema con esas caracteristicas, en este caso, el sistema propone el uso de dos enfoques:

El primer enfoque consiste en la implementacion deTF-IDF (Term Frequency–Inverse Document Frequency) y similitud del coseno, el cual permite poder medir cual es la similutd del texto dado por el usuario en comparacion con la informacion disponible en el catalogo de peliculas,para de esta manera poder determinar cual es la pelicula que mas se ajusta a la pregunta realizada por el usuario mediante la coincidencia de las palabras, por otro lado.

El segundo enfoque consiste en usar un modelo preentrenado de embeddings, este modelo permite poder representar el texto dado por el usuario y la informacion disponible en el catalogo de peliculas en un espacio vectorial, este metodo tambien usa la similitud del coseno para pode medir la similitud entre los textos y de esta manera, poder identificar las peliculas que mas se relacionen a lo que pidio el usuario. Adicionalmente, cuenta con la incorporacion de Qween, un modelo de lenguaje LLM (Large Language Model) el cual permite poder generar embeddings de manera mas eficiente y precisa, este se incluyo ya que permite poder redactar una respuesta mas completa apartir de las peliculas encontradas.   

La evaluacion se centrara en la comparacion entre estos dos enfoques para asi poder determinar la capacidad de cada enfoque para poder recuperar peliculas relevantes a partir de la pregunta realizada por el usuario, ademas de poder determinar cual de los dos enfoques es mas eficiente en cuanto a tiempo de respuesta y consumo de recursos. Las metricas de evaluacion implementadas consisten en Top-1, Recall@5, rango recíproco medio (MRR) y tiempo de respuesta.

\subsection{Objetivos}
Desarrollar un asistente conversacional que permita a los usuarios realizar consultas en español y recibir recomendaciones de películas mediante la implementacion y comparacion de dos metodos de recuperacion de informacion, uno basado en TF-IDF y otro con embeddings multilingues, ambos con similitud del coseno.

Los objetivos especi­ficos son:
\begin{enumerate}
\item Preparar un conjunto de datos de peli­culas que incluya informaciión textual relevante, como generos y sinopsis, para utilizarlo en la recuperacio­n de contenidos a partir de consultas en español.
\item Implementar un metodo de recuperacion basado en TF-IDF y similitud del coseno que permita comparar las consultas de los usuarios con la informacion textual de las peli­culas y recuperar los resultados mas relacionados.
\item Implementar un metodo de recuperacion semantica mediante embeddings multilingues y similitud del coseno, con el fin de identificar peli­culas relacionadas con la intencion de las consultas, incluso cuando existan diferencias entre los terminos utilizados por el usuario y los textos del catalogo.
\item Evaluar y comparar los dos metodos de recuperacion mediante las metricas Top-1, Recall@5, rango reciproco medio (MRR) y tiempo de respuesta, utilizando un conjunto de consultas de prueba y criterios de evaluacion definidos previamente.
\end{enumerate}

% secciones/01_introduccion/justificacion.tex
\subsection{Justificación}

Para los usuarios, el realizar una busqueda de peliculas puede resultar una tarea tediosa e incluso compleja cuando los usuarios no conocen el titulo de la pelicula o en general, no saben que pelicula quieren ver y simplemente necesitan recomendaciones de un genero en especifico. Para este tipo de consultas, es importante considerar tecnicas de procesamiento de lenguaje natural que permiten poder relacionar las necesidades expresadas en las preguntas del usuario con las caracteristicas de las peliculas en el catalogo.

Este proyecto se justifica porque permite comparar dos enfoques distintos para poder cumplir con el proposito del proyecto, TF-IDF y embeddings multilingues, ambos combinados con la similitud del coseno. Esta combinacion permite determinar si las diferencias en los modelos se reflejan en el desempeño del sistema. La eleccion de las tecnicas tambien se justifica ya que permite construir un prototipo funcional con herramientas de acceso libre y gratuito sin depender de servicios que requieran de una suscripcion o pago. 

Por otro lado, la incorporacion de un LLM permite transformar los resultados obtenidos en una respuesta mas completa y coherente para el usuario, esta capa se enfoca en la facilitacion de la presentacion de la informacion en pantalla, su funcionalidad se centra en unicamente redactar las respuestas a partir de la informacion obtenida en la busqueda, para de esta manera evitar el riesgo de que el modelo genere afirmaciones que no se encuentren respaldadas.

En general, el valor academico de este proyecto radica en que integra tecnicas de recuperacion de informacion y procesamiento de lenguaje, junto con sus metricas de evaluacion y el uso de Telegram como plataforma de interaccion con el usuario. La obtencion de resultados permite identificar las fortalezas y debilidades de cada enfoque sin asumir que una tecnica sera mejor que le otra.

De esta manera, al combinar todos estos factores se puede construir un prototipo funcional que permita a los usuarios realizar preguntas y recibir respuestas en español sobre peliculas.  
